\section{Constrained co-design problem and algorithms}\label{sec:problem}

\subsection{Workflow}
Fig.~\ref{fig:workflow} summarizes the protocol. Offline, a candidate $\xi$ is
mapped by \eqref{eq:care} and \eqref{eq:map} to a 2-DOF PIDF; the averaged
evaluator runs six scenarios and the switched screen runs three short steady
windows; Deb ordering ranks the record. After 240 global and 60 local calls the
best record of each run is frozen, and only then replayed on the full switched
mission, which is never visible to the search.

\begin{figure*}[t]
\centering
\begin{tikzpicture}[
  font=\footnotesize, node distance=4mm and 5mm,
  box/.style={draw=black!60, rounded corners=2pt, align=center, inner sep=3pt, minimum height=9mm, fill=white},
  hl/.style={box, fill=blue!6},
  sw/.style={box, fill=orange!8},
  arr/.style={-{Stealth[length=2mm]}, thick, black!70},
  darr/.style={arr, dashed}]
\node[box, text width=28mm] (plant) {Non-ideal Buck\\ Eqs.~\eqref{eq:node}--\eqref{eq:dff}, $G_{vd}(s)$};
\node[hl, right=of plant, text width=30mm] (lqi) {LQI weights $\xi$\\ Bryson $Q,R_d$ $\to$ CARE \eqref{eq:care}};
\node[hl, right=of lqi, text width=30mm] (map) {Exact map \eqref{eq:map}\\ $K\to(K_p,K_i,K_d,N_C,1,0,K_b)$};
\node[box, right=of map, text width=30mm] (avg) {Averaged evaluator\\ 6 scenarios, $J_s$, $g_1..g_{36}$};
\node[sw, below=of avg, text width=30mm] (scr) {Switched screen\\ 24$\to$1$\to$46$\to$24 V, 34 ms: $g_{37}..g_{39}$};
\node[box, left=of scr, text width=30mm] (deb) {Deb record\\ feasible-first, no penalty};
\node[hl, left=of deb, text width=30mm] (opt) {7 algorithms, 30 seeds\\ 240 global + 60 refineBO};
\node[box, left=of opt, text width=28mm] (cert) {Certificates\\ LQR margins, D-region LMI, Prop.~\ref{prop:struct}};
\node[box, below=of deb, text width=30mm] (freeze) {Freeze $\theta^\star$\\ (no switched mission data seen)};
\node[sw, right=of freeze, text width=30mm] (val) {Switched mission 324 ms\\ energy balance, specs, comparators};
\draw[arr] (plant) -- (lqi);
\draw[arr] (lqi) -- (map);
\draw[arr] (map) -- (avg);
\draw[arr] (avg) -- (scr);
\draw[arr] (scr) -- (deb);
\draw[arr] (deb) -- (opt);
\draw[arr] (opt.north) |- ++(0,3mm) -| (lqi.south);
\draw[arr] (lqi.south west) ++(2mm,0) |- (cert.north);
\draw[arr] (opt.south) |- (freeze.west);
\draw[darr] (freeze) -- (val);
\end{tikzpicture}

\caption{Formula-first CMO--LQI--PIDF workflow. Solid arrows: evaluation inside the search; dashed arrow: acceptance replay after the freeze. The switched screen is the only switched information used during the search.}
\label{fig:workflow}
\end{figure*}

\subsection{Scenarios and objective}
The averaged evaluator integrates \eqref{eq:avg} with the hybrid guard $i_L\ge0$
(RK4, \SI{2}{\micro s}) under the sampled control layer of
Section~\ref{sec:theory}. Six scenarios are declared (Table~\ref{tab:scen}); each
starts from the analytical equilibrium preceding its first event. For scenario
$s$ the seven criteria are the loss intensity $E_{loss}/E_u$, IAE, ISE, global
ITAE, event-reset ITAE, RMSE and MAE, each divided by its value for the
normalization record $\xi_0$. The scalar objective is the mean of the $6\times7$
ratios,
\begin{equation}
J_s(\xi)=\frac{1}{42}\sum_{s=1}^{6}\mathbf 1^\top r_s(\xi),\qquad J_s(\xi_0)=1 .
\label{eq:J}
\end{equation}
MAE is IAE divided by the horizon and RMSE is the square root of ISE divided by
it, so two of the seven entries are redundant; the induced weights (2/7 on the
IAE family, 2/7 on the ISE family, 1/7 on each other criterion) are declared
rather than left implicit. The objective ranks feasible candidates only.

\begin{table}[t]
\centering
\caption{Averaged scenarios (the source is \SI{48}{V} in all of them).}
\label{tab:scen}
\footnotesize
\begin{tabular}{@{}llll@{}}
\toprule
ID & $T$ (ms) & reference & load pulse (+0.5 A)\\
\midrule
S0 & 90 & 24$\to$1$\to$46 V at 30, 60 ms & 15--25 ms\\
S1a & 40 & 24$\to$1 V at 10 ms & --\\
S1b & 40 & 1$\to$46 V at 10 ms & --\\
S1c & 40 & 46$\to$24 V at 10 ms & --\\
S2a & 40 & 24 V & 10--20 ms\\
S2b & 40 & 46 V & 10--20 ms\\
\bottomrule
\end{tabular}
\end{table}

\subsection{Thirty-nine inequalities in engineering units}
For each scenario six inequalities $g_{s,j}\le0$ are evaluated separately:
\begin{equation}
\begin{aligned}
g_{s,1}&=\frac{\max_t\,\sigma_k\big(v_o-r_k\big)}{\SI{1.5}{V}}-1 &&\text{overshoot beyond each new level},\\
g_{s,2}&=\frac{\max_{t\in\mathcal T_p}|r-v_o|}{\SI{0.5}{V}}-1 &&\text{load-pulse deviation},\\
g_{s,3}&=\frac{\max_t i_L+\Delta I_{L,pp}^{\max}/2}{\SI{23}{A}}-1 &&\text{device current},\\
g_{s,4}&=\frac{\max_k|\bar e_k|}{\SI{25}{mV}}-1 &&\text{steady error},\\
g_{s,5}&=\frac{|E_u/E_{u,0}-1|}{0.005}-1 &&\text{same useful task},\\
g_{s,6}&=\frac{\rho_{\ell,s}-1}{0.005}-1,\;\; \rho_{\ell,s}=\frac{E_{loss}/E_u}{(E_{loss}/E_u)_{\xi_0}} &&\text{loss intensity},
\end{aligned}
\label{eq:g}
\end{equation}
where $\sigma_k$ is the sign of the $k$-th reference step, $\mathcal T_p$ runs from
the onset of the load pulse to the next event, $\bar e_k$ is the mean error over
the last \SI{2}{ms} of each constant segment and \SI{23}{A} is the IRF540N
continuous drain rating at \SI{100}{\celsius}. The averaged peak current is
augmented by the worst half ripple $\Delta I_{L,pp}^{\max}/2=\SI{1.17}{A}$
(Table~\ref{tab:tf}). A family that does not apply to a scenario (no step, no
load pulse) is set to $-1$. These 36 inequalities are completed by three measured ones,
\begin{equation}
g_{37}=\frac{\max_{k}\,\mathrm{pp}(d)}{0.20}-1,\quad
g_{38}=\frac{\max_{k}|\bar e_{sw,k}|}{\SI{25}{mV}}-1,\quad
g_{39}=\frac{\max_t i_{L,sw}}{\SI{23}{A}}-1,
\label{eq:gsw}
\end{equation}
obtained from one run of the switched replica over the transition sequence
24$\to$1$\to$46$\to$24~V (steps at 2, 14 and \SI{24}{ms}, \SI{34}{ms} in total,
\SI{1}{\micro s} RK4 with exact PWM edges). The duty peak-to-peak and the mean
error are measured over the last \SI{4}{ms} of each level. A candidate is
feasible when $\sum_j\max(g_j,0)\le10^{-4}$; the screen costs about \SI{15}{ms}
per candidate against \SI{14}{ms} for the six averaged scenarios.

\paragraph{Why the screen contains large transitions}
The protocol was frozen after two pilot campaigns that are reported, not
hidden. Pilot~1 used the analytical ripple bound of the earlier version of this
study, $\rho_{rip}=\max_V|C_{fb}(z{=}{-1})|\Delta V_{o,pp}(V)/2\le0.20$, where
$C_{fb}$ is the discrete feedback path evaluated at the controller Nyquist
frequency, which coincides with the carrier frequency because $T_s=1/(2f_s)$.
Every branch winner it produced satisfied the bound ($\rho_{rip}\approx0.07$--$0.19$)
and yet, on the switched model, several sustained a limit cycle with duty swings
of 0.7--0.9 at \SI{24}{V} and offsets of 0.3--\SI{1.8}{V}. Pilot~2 replaced the
bound by steady \SI{10}{ms} switched windows started \emph{from equilibrium} at
1, 24 and \SI{46}{V}. Its best design (178 of 180 runs feasible) passed that
screen, but on the reference mission it entered, after the large
24$\to$1~V transition, a different limit cycle in which one sample in four
clips at $d_{\min}$, leaving a \SI{-42}{mV} offset at \SI{1}{V}. The loop thus
holds \emph{coexisting} switching limit cycles, and the one reached depends on
the transient that precedes it; this is why the final screen is a transition
sequence. $\rho_{rip}$ is kept as a reported diagnostic only.

Feasibility is stated in absolute units, not as non-inferiority against
$\xi_0$. With ratio gates, 24 of 36 constraints are active at $\xi_0$ and a
gentler normalization record leaves no candidate that dominates it; with
absolute gates every controller family of Section~\ref{sec:families} is judged
against the same specification, and $\xi_0$ itself is infeasible because it
overshoots by \SI{6.2}{V} on the 1$\to$46~V step.

\subsection{Algorithms and budgets}\label{sec:algos}
Constraint handling uses Deb's rules \citep{Deb2000}: feasible beats
infeasible, feasible records are ranked by $J_s$, infeasible ones by aggregate
violation; no penalty weight is tuned. All methods work in the normalized box
$u\in[0,1]^n$, receive the matched seed $2026091200+k$, $k=1,\dots,30$, and the same
initial population of 12: $\xi_0$ first, six further individuals within a
normalized radius 0.05 of it, five uniform. One counted call is one complete
evaluation (six averaged scenarios, the switched screen, 39 inequalities).
Table~\ref{tab:algos} lists the operators.

\begin{table*}[t]
\centering
\caption{Algorithms 1--7. Population 12, 240 global calls, then 60 refineBO calls initialized from the same run's history (300 calls per run).}
\label{tab:algos}
\footnotesize
\begin{tabular}{@{}lp{6.0cm}p{4.3cm}p{4.4cm}@{}}
\toprule
Alg. & Exploration / update & Exploitation / survival & Fixed parameters\\
\midrule
1 PSO \citep{Kennedy1995} & velocity with independent cognitive and social draws & personal and global Deb incumbents & $w=0.9-0.5\tau$, $c_1=c_2=1.7$, $|v|\le0.2$\\
2 GA & binary Deb tournaments, SBX crossover & polynomial mutation, elitist $(\mu+\lambda)$ & $\eta_c=15$, $\eta_m=20$, $p_m=1/n$\\
3 ABC \citep{Karaboga2008} & one-coordinate differential move (employed, onlooker) & greedy Deb replacement, scout & limit $\max(5,\mathrm{round}(0.5\,n_pn))$\\
4 pAEABC & ABC move scaled by diversity $s_k=\mathrm{clip}(0.5+0.15/D_k,0.5,2)$, extra coordinate kick w.p.\ $0.25(1-\tau)$ & as ABC & inspired by \citet{Alsamia2024}\\
5 pIGWO & three-leader GWO with $a=2(1-\tau^2)$, Gaussian kick w.p.\ $0.15(1-\tau)+0.02$ & elitist pooling & inspired by \citet{Hou2022}\\
6 pIGWO--DLH & linear $a=2-2k/K$ GWO trial and an independently evaluated radial local-hunt trial & better trial replaces the wolf (non-elitist) & $R_i=\|x_i-x_\alpha\|/2$; inspired by \citet{NadimiShahraki2021}\\
7 refineBO & shared-kernel Matérn-5/2 GPs on $J_s$ and $g_1,\dots,g_{39}$; pool of 1000 around the 10 Deb-best records, radius $0.04\to0.003$ & $\alpha=P_f(\mathrm{EI}+0.02\sigma_J)+10^{-4}P_f$, $P_f=\prod_j\Phi(-\mu_j/\sigma_j)$; restoration mode if no feasible incumbent & 60 calls\\
\bottomrule
\end{tabular}
\end{table*}

The prefix ``p'' marks project variants that are inspired by, but not claimed to
be operator-identical with, the cited algorithms. $\tau$ is the fraction of the
budget consumed. The product $P_f$ of marginal feasibility probabilities is an
independence approximation used for ordering only; every accepted point is
checked by the exact evaluator. Wall time is logged but never interpreted as an
energy or performance quantity.

\subsection{Controller families and fixed-rule comparators}\label{sec:families}
To separate the value of the parameterization from that of the optimizer, the
identical protocol (evaluator, 39 inequalities, 7 algorithms, 30 seeds,
300 calls) is run on six families:
\begin{itemize}
\item \textbf{CMO--LQI--PIDF}: $\xi$ of \eqref{eq:xi} $\to$ \eqref{eq:care}
  $\to$ \eqref{eq:map}; 4 variables, reads $v_o$ only;
\item \textbf{CMO--LQI--SF}: the same $\xi$ and the same continuous gain, applied
  as full-state LQI with measured $i_L$ and $i_C$ (the twin of
  Proposition~\ref{prop:struct});
\item \textbf{PIDF-direct}: the same 2-DOF structure ($b=1$, $c=0$, $N=N_C$) with
  $(\log K_p,\log K_i,\log K_d,\log K_b)$ searched directly over the gain ranges
  spanned by the LQI box;
\item \textbf{PIDF7-direct}: $K_p,K_i,K_d,N,b,c,K_b$ all free (7 variables);
\item \textbf{PID-direct}: backward-difference derivative on $-v_o$, $b=1$,
  $(K_p,K_i,K_d,K_b)$ free;
\item \textbf{PI-direct}: $(K_p,K_i,K_b)$ free.
\end{itemize}

\paragraph{Fixed-rule comparators}\label{sec:comparators}
Six classical designs share the control layer and are synthesized at \SI{24}{V},
$T_s=\SI{10}{\micro s}$, with the delay $1.5T_s$ (computation plus ZOH) in every
frequency-domain rule. \emph{PI}: crossover $0.03f_s$, phase margin
\SI{60}{\degree}, $K_b=K_i/K_p$. \emph{PID}: zeros cancelling the $LC$ pair of
$\Delta(s)$, crossover set for \SI{60}{\degree} of phase margin including the ESR
zero, backward-difference derivative on $-v_o$, $b=1$, $c=0$,
$K_b=1/\sqrt{T_iT_d}$. \emph{PIDF}: the Type-III rule, i.e.\ the PID zeros plus a
filter pole on the ESR zero $N=N_C$, \SI{60}{\degree} of phase margin.
\emph{LQR}: discrete Riccati solution on the ZOH model with the Bryson limits of
$\xi_0$, measured $i_L$ and $v_C=v_o-r_Ci_C$, no integral action. \emph{LQG}: the
LQR gain with a steady-state Kalman predictor on $v_o$ only (process noise
0.01 in duty, measurement noise \SI{0.05}{V}). \emph{LQI}: discrete Riccati
solution of the integral-augmented ZOH model with the same weights, measured
state, $K_b=\omega_B$. None of them is tuned by the proposed search; they
represent what the standard rules deliver on this plant.

\subsection{Statistics and validation}
Methods are compared as matched blocks over the 30 seeds with the Friedman test
and Holm-adjusted paired Wilcoxon signed-rank tests \citep{Demsar2006}; a run's
score is $J_s$ if feasible and $10^3$ plus its violation otherwise, which is
consistent with Deb ordering. Families are compared pairwise on the same
(method, seed) cells. The retained controller of each family is the best
feasible record over its 180 runs; it is frozen and replayed on the
\SI{324}{ms} reference mission of the Simulink model
(24$\to$1$\to$46~V at 108 and 216~ms, +0.5~A load pulse at 54--90~ms, initial
state at the \SI{24}{V} equilibrium), with the energy balance of
\eqref{eq:energy}.
